\chapter{Architecture technique de l'application}
\label{ch:archi}

\section{Introduction}

Ce chapitre décrit l'architecture logicielle de \nomapplication\ : organisation du back-end Java, structure du front-end Thymeleaf, et flux de données entre les couches.

\section{Architecture globale}

L'application suit une \textbf{architecture en couches} (layered monolith) :

\begin{center}
\begin{tabular}{c}
    \textbf{Navigateur (HTML/CSS/JS)} \\
    $\downarrow$ \\
    \textbf{Contrôleurs} (\texttt{web.controller}) \\
    $\downarrow$ \\
    \textbf{Services métier} (\texttt{service}) — \texttt{@Transactional} \\
    $\downarrow$ \\
    \textbf{Repositories} (\texttt{repository}) — Spring Data JPA \\
    $\downarrow$ \\
    \textbf{PostgreSQL} (Flyway V1--V11) \\
\end{tabular}
\end{center}

Transversal : \textbf{Spring Security} (filtre HTTP), \textbf{AuditService}, \textbf{NotificationModelAdvice}.

\section{Arborescence du back-end}

Le code source se trouve dans \texttt{src/main/java/com/banque/agence/}.

\subsection{Point d'entrée et configuration}

\begin{table}[H]
    \centering
    \caption{Packages de configuration}
    \label{tab:config}
    \small
    \begin{tabular}{@{}lp{9cm}@{}}
        \toprule
        \textbf{Fichier / classe} & \textbf{Rôle} \\
        \midrule
        \texttt{BanqueAgenceApplication.java} & Point d'entrée Spring Boot \\
        \texttt{config/SecurityConfig.java} & Chaîne de sécurité, rôles, redirection login \\
        \texttt{config/DevUserInitializer.java} & Seed utilisateurs (admin, agent, chef) \\
        \texttt{config/DevDemoDataInitializer.java} & Seed données métier (clients, comptes…) \\
        \texttt{config/DemoPortalSync.java} & Activation portail client démo \\
        \texttt{application.yml} & Port 8081, profil actif \\
        \texttt{application-dev.yml} & Connexion PostgreSQL, flags seed \\
        \bottomrule
    \end{tabular}
\end{table}

\subsection{Couche domaine}

\begin{table}[H]
    \centering
    \caption{Entités et énumérations}
    \label{tab:domain}
    \small
    \begin{tabular}{@{}lp{9cm}@{}}
        \toprule
        \textbf{Package} & \textbf{Contenu} \\
        \midrule
        \texttt{domain/entity/} & \texttt{User}, \texttt{Client}, \texttt{Account}, \texttt{Transaction}, \texttt{BillProvider}, \texttt{BillPayment}, \texttt{CheckbookOrder}, \texttt{Notification}, \texttt{AuditLog} \\
        \texttt{domain/enums/} & \texttt{UserRole}, \texttt{AccountType}, \texttt{AccountStatus}, \texttt{TransactionType}, \texttt{ClientStatus}, \texttt{CheckbookOrderStatus}, \texttt{NotificationType}… \\
        \bottomrule
    \end{tabular}
\end{table}

\subsection{Couche persistance}

\begin{itemize}[leftmargin=*]
    \item \texttt{repository/} : interfaces Spring Data JPA (\texttt{ClientRepository}, \texttt{AccountRepository}, \texttt{TransactionRepository}…).
    \item \texttt{TransactionSpecifications.java} : filtres dynamiques pour l'historique.
    \item Migrations SQL : \texttt{src/main/resources/db/migration/V1\_\_} à \texttt{V11\_\_}.
\end{itemize}

\subsection{Couche métier (services)}

\begin{table}[H]
    \centering
    \caption{Services principaux}
    \label{tab:services}
    \small
    \begin{tabular}{@{}lp{9cm}@{}}
        \toprule
        \textbf{Service} & \textbf{Responsabilité} \\
        \midrule
        \texttt{ClientService} & CRUD clients, recherche, changement statut \\
        \texttt{AccountService} & Ouverture, blocage, clôture comptes \\
        \texttt{TransactionService} & Dépôt, retrait, virement (\texttt{@Transactional}) \\
        \texttt{BillPaymentService} & Paiement facture + transaction liée \\
        \texttt{CheckbookOrderService} & Workflow chéquier sans impact solde \\
        \texttt{DashboardService} & KPIs et opérations récentes \\
        \texttt{NotificationService} & Création et lecture notifications \\
        \texttt{ClientPortalService} & Données portail client (lecture seule) \\
        \texttt{ReceiptPdfService} & Génération PDF OpenPDF \\
        \texttt{UserService} & Administration personnel \\
        \texttt{AuditService} & Journalisation actions sensibles \\
        \bottomrule
    \end{tabular}
\end{table}

\subsection{Couche web (contrôleurs)}

\textbf{Agence} (\texttt{web.controller/}) :
\begin{itemize}[leftmargin=*]
    \item \texttt{HomeController} — landing \texttt{/}
    \item \texttt{LoginController}, \texttt{DashboardController}
    \item \texttt{ClientController}, \texttt{AccountController}
    \item \texttt{OperationController} — dépôt, retrait, virement, facture
    \item \texttt{TransactionController} — historique, reçu HTML/PDF
    \item \texttt{CheckbookOrderController}, \texttt{NotificationController}
    \item \texttt{AdminUserController}, \texttt{AuditController}
\end{itemize}

\textbf{Portail client} (\texttt{web.controller.portal/}) :
\begin{itemize}[leftmargin=*]
    \item \texttt{PortalDashboardController}, \texttt{PortalAccountController}
    \item \texttt{PortalTransactionController}, \texttt{PortalCheckbookController}
    \item \texttt{PortalNotificationController}
\end{itemize}

\textbf{DTO} (\texttt{web.dto/}) : objets formulaire (\texttt{ClientForm}, \texttt{OperationForm}, \texttt{BillPaymentForm}…).

\subsection{Sécurité}

\begin{itemize}[leftmargin=*]
    \item \texttt{UserDetailsServiceImpl} : charge \texttt{User} (personnel) ou \texttt{Client} (portail).
    \item \texttt{SecurityUser} / \texttt{SecurityClient} : wrappers Spring Security.
    \item Routes : \texttt{/portal/**} $\rightarrow$ ROLE\_CLIENT ; \texttt{/admin/**} $\rightarrow$ ROLE\_ADMIN.
\end{itemize}

\section{Front-end (Thymeleaf)}

\subsection{Identité visuelle et charte graphique}

L'interface a été entièrement repensée autour de l'identité \textbf{AmanaBank}. La palette repose sur trois familles de couleurs définies en variables CSS (\texttt{:root}) :

\begin{itemize}[leftmargin=*]
    \item \textbf{Marron} (\texttt{\#4B3621} à \texttt{\#7A6352}) : barre latérale, textes principaux ;
    \item \textbf{Beige} (\texttt{\#FDFBF8} à \texttt{\#EDE6DA}) : arrière-plans, cartes et en-têtes ;
    \item \textbf{Or} (\texttt{\#A6893A}) : boutons d'action, liens actifs et accents.
\end{itemize}

La police \textbf{DM Sans} (Google Fonts) assure une lecture confortable sur tous les écrans. Le logo officiel (\texttt{amana-logo.png} pour la vitrine, \texttt{AmanaBank.png} dans les espaces connectés) et le favicon renforcent la cohérence de marque.

Deux feuilles de style complémentaires structurent le rendu :
\begin{itemize}[leftmargin=*]
    \item \texttt{landing.css} — page publique autonome (header fixe, hero plein écran, grilles de services) ;
    \item \texttt{style.css} — espaces connectés agence et portail (sidebar, cartes KPI, formulaires, tableaux).
\end{itemize}

\subsection{Structure des templates}

Les vues se trouvent dans \texttt{src/main/resources/templates/}.

\begin{table}[H]
    \centering
    \caption{Organisation des templates}
    \label{tab:templates}
    \small
    \begin{tabular}{@{}lp{9cm}@{}}
        \toprule
        \textbf{Dossier / fichier} & \textbf{Rôle} \\
        \midrule
        \texttt{landing.html} & Page vitrine publique AmanaBank \\
        \texttt{login.html} & Connexion unique personnel + client \\
        \texttt{layout/main.html} & Layout agence (sidebar, cloche, menu) \\
        \texttt{dashboard.html} & Tableau de bord KPIs et alertes \\
        \texttt{clients/} & Liste, formulaire, fiche détail client \\
        \texttt{accounts/} & Liste, ouverture, détail, relevé \\
        \texttt{operations/} & Dépôt, retrait, virement, paiement facture \\
        \texttt{transactions/} & Historique, reçu imprimable \\
        \texttt{checkbook-orders/} & Liste, demande, détail workflow \\
        \texttt{notifications/} & Centre de notifications personnel \\
        \texttt{admin/users/} & Gestion utilisateurs internes \\
        \texttt{admin/audit/} & Journal d'audit \\
        \texttt{portal/layout.html} & Layout espace client \\
        \texttt{portal/} & Dashboard, comptes, historique, chéquiers, notifications \\
        \texttt{error/403.html}, \texttt{404.html} & Pages d'erreur \\
        \bottomrule
    \end{tabular}
\end{table}

\subsection{Ressources statiques}

\begin{itemize}[leftmargin=*]
    \item \texttt{static/css/style.css} — thème unifié agence et portail (sidebar, cartes KPI, login).
    \item \texttt{static/css/landing.css} — vitrine marketing (hero, statistiques, services, footer).
    \item \texttt{static/js/landing-header.js} — en-tête fixe avec masquage au défilement.
    \item \texttt{static/images/} — logos, favicon, visuels (\texttt{hero.jpg}, \texttt{branch.jpg}, \texttt{digital.jpg}…).
\end{itemize}

\subsection{Page d'accueil publique (\texttt{landing.html})}

La landing est une page \textbf{autonome} (sans layout Thymeleaf partagé) structurée en sections ancrées :

\begin{enumerate}[leftmargin=*]
    \item \textbf{En-tête fixe} — logo, navigation (Services, Espace digital, Engagements, Contact), boutons Connexion et Ouvrir un compte.
    \item \textbf{Hero} — image de fond (\texttt{hero.jpg}), accroche institutionnelle et appels à l'action.
    \item \textbf{Chiffres clés} — bandeau de statistiques (expérience, clients, agences, satisfaction).
    \item \textbf{Services} — trois cartes illustrées (comptes, financement, services digitaux).
    \item \textbf{Espace client} — présentation du portail avec liste à puces et visuel mobile.
    \item \textbf{Engagements} — valeurs de la banque (sécurité, proximité, transparence).
    \item \textbf{Contact} — appel à l'action et coordonnées.
    \item \textbf{Pied de page} — liens, adresse et mentions légales.
\end{enumerate}

Le contrôleur \texttt{HomeController} sert cette page sur la route \texttt{/} sans authentification.

\subsection{Connexion unifiée (\texttt{login.html})}

Un seul écran de connexion accueille le \textbf{personnel} (username) et les \textbf{clients} (CIN ou numéro client). Le formulaire POST vers \texttt{/login} est géré par Spring Security ; \texttt{UserDetailsServiceImpl} détermine le type de compte chargé. L'interface affiche des encadrés d'aide distincts pour chaque profil et un lien de retour vers la landing.

\subsection{Layouts connectés}

\textbf{Agence} (\texttt{layout/main.html}) : sidebar marron avec navigation par rôle (\texttt{sec:authorize}), en-tête de page, cloche de notifications pour le chef d'agence, zone de contenu sur fond beige.

\textbf{Portail client} (\texttt{portal/layout.html}) : même structure visuelle avec badge «\,Espace client\,», menu réduit (comptes, historique, chéquiers) et cloche de notifications personnelles.

Les tableaux de bord (\texttt{dashboard.html}, \texttt{portal/dashboard.html}) réutilisent le composant \texttt{stat-card--kpi} : indicateur chiffré, icône colorée, lien d'action rapide vers le module concerné.

\subsection{Composants transverses}

\begin{itemize}[leftmargin=*]
    \item \texttt{NotificationModelAdvice} : injecte \texttt{unreadNotificationCount} dans toutes les vues.
    \item \texttt{thymeleaf-extras-springsecurity6} : affichage conditionnel selon rôle (\texttt{sec:authorize}).
    \item Messages flash (\texttt{RedirectAttributes}) pour succès et erreurs métier.
\end{itemize}

\section{Conclusion}

Ce chapitre a présenté l'architecture en couches de \nomapplication, l'organisation du back-end Java et la structure du front-end Thymeleaf. La refonte de l'interface — charte AmanaBank, landing autonome, layouts à sidebar et composants KPI réutilisables — assure une cohérence visuelle entre l'espace agence et le portail client. Le chapitre suivant illustre ces interfaces par des captures d'écran commentées.
